% Source: https://artofproblemsolving.com/community/c3h1441503_2013_national_sprint_round_16
% Stable identity: aops-topic-1441503-q16
The regions have areas $\pi$, $4\pi-\pi=3\pi$ and $9\pi-4\pi=5\pi$ out of $9\pi$, so one dart scores $10$, $5$, $1$ with probabilities $\dfrac{1}{9}$, $\dfrac{3}{9}$, $\dfrac{5}{9}$.

\begin{asy}
usepackage("amsmath");
size(4.5cm);
defaultpen(fontsize(10pt));
fill(circle((0,0),3),paleyellow); fill(circle((0,0),2),lightcyan); fill(circle((0,0),1),pink);
draw(circle((0,0),3)); draw(circle((0,0),2)); draw(circle((0,0),1));
label("$10$",(0,0.45)); label("$5$",(0,1.5)); label("$1$",(0,2.5));
draw((0,0)--(3,0),gray); dot((0,0)); dot((1,0),gray); dot((2,0),gray); dot((3,0),gray);
label("$1$",(1,0),SE,gray+fontsize(8pt)); label("$2$",(2,0),SE,gray+fontsize(8pt)); label("$3$",(3,0),SE,gray+fontsize(8pt));
label("\mbox{}",(min(currentpicture,user=true)+max(currentpicture,user=true))/2,black); // keeps the text after the figure black when LaTeX typesets the labels inline
\end{asy}

Modulo $3$ the scores are $10\equiv1$, $5\equiv2$, $1\equiv1$. So each dart is $\equiv1$ with probability $\dfrac{1}{9}+\dfrac{5}{9}=\dfrac{2}{3}$ and $\equiv2$ with probability $\dfrac{1}{3}$. A sum of three residues each equal to $1$ or $2$ is divisible by $3$ only if all three are equal ($1+1+1$ or $2+2+2$; a mixture gives $4$ or $5$). Hence
\[
P=\left(\dfrac{2}{3}\right)^3+\left(\dfrac{1}{3}\right)^3=\dfrac{8+1}{27}=\boxed{\dfrac{1}{3}}.
\]
The ``shortcut'' is exactly this reduction modulo $3$: only two residue classes occur.
